% Example: Theorem, lemma, and proof environments with amsthm
% Demonstrates declaring theorem-like environments with \newtheorem
% and writing structured proofs using the proof environment.
\documentclass[12pt, a4paper]{article}
\usepackage{amsmath, amssymb, amsthm}
\usepackage{fontspec}

% Declare theorem-like environments
\newtheorem{theorem}{Theorem}
\newtheorem{lemma}[theorem]{Lemma}
\newtheorem{corollary}[theorem]{Corollary}

\begin{document}

\section*{Theorems and Proofs}

\begin{theorem}[Pythagoras]
In a right-angled triangle, the square of the hypotenuse equals the
sum of the squares of the other two sides:
\[
  a^2 + b^2 = c^2.
\]
\end{theorem}

\begin{proof}
Consider a right-angled triangle with legs \(a\) and \(b\) and
hypotenuse \(c\). By constructing a square on each side and comparing
areas, one obtains \(a^2 + b^2 = c^2\).
\end{proof}

\begin{lemma}
For any real numbers \(x\) and \(y\),
\[
  (x + y)^2 = x^2 + 2xy + y^2.
\]
\end{lemma}

\begin{proof}
Expanding the left-hand side:
\begin{align*}
  (x + y)^2 &= (x + y)(x + y) \\
            &= x^2 + xy + yx + y^2 \\
            &= x^2 + 2xy + y^2. \qedhere
\end{align*}
\end{proof}

\begin{corollary}
For any real number \(x\),
\[
  (x + 1)^2 = x^2 + 2x + 1.
\]
\end{corollary}

\begin{proof}
Substitute \(y = 1\) into the lemma.
\end{proof}

\end{document}
